\section{Purpose}


As part of maintaining IANZ accreditation to the standard ISO/IEC 17025:2017, MSL is required to proactively seek feedback from its customers.  Specifically, clause 8.6.2 of that standard states ‘The laboratory shall seek feedback both positive and negative, from its customers.  The feedback shall be analysed to improve the management system, laboratory activities and customer service.’

These guidelines provide practical guidance for conducting and recording client call-backs.

\section{Selecting clients}

The MSL Quality Manager will create and maintain a list of key clients and assign call backs.  Typically, this will be New Zealand based calibration laboratories but may include others too.  The MSL Director will follow up on any overdue call backs.

The frequency of call back will generally be 2 yearly and where possible will be timed to coincide with recent work.  However some clients may warrant a shorter interval, and those for whom MSL has done no work in the last 2 years can be excluded.

\section{Conducting the call-back}
\label{s:conducting_call_backs}

The phone calls will be in the role of a client relationship manager.  That is, each client on the call back list will have an assigned team manager who will be the default person to do the call back each time they are due.  This assignee may be changed as circumstances require.  Calls will be semi-structured with possible topics described below but the onus is on the person making the call to adapt the conversation with what seems relevant at the time.

Overall, how have the recent jobs been with MSL?
\begin{itemize}
\item How would you describe the quality and ease of communication with MSL?  This can include any outstanding enquiries and quote stages as well, not just jobs completed.
\item Asking about the clients upcoming challenges or any new services they are offering. This could lead to discussion on how MSL may be able to help, or if there are any new services that MSL should consider offering.
\item Knowing a little more about their operations could help understand the need for new measurement services to be developed at MSL (or inform which of MSL services are most critical to them), or indirectly inform the potential to apply cost increases/value-based pricing.
\item Essentially try and understand what we could or need to do better.
\item Promoting other potential services from MSL, not necessarily fee based.  For example, technical guides and training courses.
\end{itemize}

\section{Recording the outcome}

At the conclusion of the phone call, form MSLF.004 containing a summary of the call will be uploaded to the MSL Quality Management System.
